% Imporditud: tools/import_eksperimendid.py
% Allikas: Eesti füüsikaolümpiaadi eksperimendi ülesannete
%   kogu (Rael Kalda, 2023), ülesanne Ü103, lahendus L103.
\setAuthor{EFO žürii}
\setRound{lõppvoor}
\setYear{2011}
\setNumber{G E1}
\setDifficulty{3}
\setTopic{vedelike mehaanika}
\setVahendid{süstal (ruumalaga V = 10 ml), anum veega, statiivile kinnitatud joonlaud, stopper}
\setPunktid{12}
\setAllikas{Eksperimendi kogu Ü103, lahendus lk 79}
\setLahendusseis{toores}

\prob{Süstal}
Määrake süstla nõela siseläbimõõt.

\hint

\solu
Tõmbame süstla vett täis, suuname selle üles ja tühjendame surudes ühtlaselt, nii

et veejuga kerkib kogu aeg enam-vähem samale kõrgusele. Mõõdame joonlaua abil

veejoa kõrguse h ja stopperiga süstla tühjenemiseks kuluva aja. Statiivi saab ka-

sutada joonlaua vertikaalsena hoidmiseks. Energia jäävusest $\pi$4vde2evjota=ja$\pi$$o_4$ kds2tsaa2mghe. veejoa kiiruseks süstlast väljumisel v = 2gh; veekulu V =

Siit saame avaldada süstla sisediameetri

d = 4 V (2gh)$-$1/4.

$\pi$t

Mõõtmistulemuste t = (54 $\pm$ 1) s ja h = (12 $\pm$ 2) cm ning väärtuse V = 10 ml põhjal saame d $\approx$ 0,36 mm.

Märkus: Tulemuse mõõtemääramatust on võimalik leida seosest

$\Delta$d $\Delta$t $\Delta$h $\Delta$V =+ + $\approx$ 0,08, d 2t 4h 2V

kus ruumala V määramatuseks on võetud $\Delta$V = 0,5 ml. Seega d = (0,36 $\pm$ 0,03) mm. Nii et vaatamata suhteliselt ebatäpsetele üksikmõõtmistele ei tule suhteline viga kokkuvõttes eriti suur.
\probend
